% ==========================================
% Kapitel 6: Komplexer Logarithmus
% ==========================================
\section{Komplexer Logarithmus}

\subsection*{Motivation: Stammfunktion des Logarithmus}
Das Ziel ist es, eine Stammfunktion zu $f(z) = \frac{1}{z}$ auf dem Gebiet $\C \setminus (-\infty, 0]$ zu finden. Da die Exponentialfunktion $\exp(z) = e^x(\cos y + \I \sin y)$ periodisch mit $2\pi\I$ ist, ist sie auf einem Streifen der Höhe $2\pi$ bijektiv.

\begin{defi}{6.1 Komplexer Logarithmus}
Sei $\theta \in \R$ fest. Die Umkehrfunktion der Exponentialfunktion
\[
    \exp: \{z \in \C : \theta < \operatorname{Im}(z) < \theta + 2\pi\} \longrightarrow \C \setminus \{e^{\I \theta}t : t \ge 0\}
\]
heißt ein \textbf{Zweig des komplexen Logarithmus}.
\end{defi}

\begin{tcolorbox}[colback=white, colframe=black!70, title={Der Hauptzweig}]
Die Umkehrfunktion zu 
\[
    \exp: \{z \in \C : -\pi < \operatorname{Im}(z) < \pi\} \longrightarrow \C \setminus (-\infty, 0]
\]
heißt \textbf{Hauptzweig des komplexen Logarithmus}, kurz \textbf{log}.
Es gilt:
\[
    \log: \C \setminus (-\infty, 0] \longrightarrow \{z \in \C : -\pi < \operatorname{Im}(z) < \pi\}
\]
Ist kein spezieller Zweig angegeben, so bezeichnet \textbf{log} stets den Hauptzweig.
\end{tcolorbox}

\begin{warnung}{Zweige des Logarithmus}
Andere Zweige werden ebenfalls oft als "log" bezeichnet, wobei dann der gewählte Bereich (Zweig) spezifiziert werden muss. Die Unstetigkeit beim "Durchschreiten" des Schnitts ($\theta$ bis $\theta + 2\pi$) verdeutlicht, warum man auf die "geschlitzte Ebene" $\C \setminus \{e^{\I \theta}t : t \ge 0\}$ angewiesen ist.
\end{warnung}

\begin{tcolorbox}[colback=white, colframe=black!70, title={Bemerkung 6.2: Eigenschaften des Hauptzweiges}]
\begin{itemize}
    \item[i)] $\log z = \ln|z| + \I \varphi$, wobei $\varphi \in (-\pi, \pi)$ das Argument von $z$ ist. Es gilt $e^{\log z} = z$ für alle $z \in \C \setminus (-\infty, 0]$.
    \item[ii)] $\log e^z = z$ für alle $z \in \C$ mit $-\pi < \operatorname{Im}(z) < \pi$.
    \item[iii)] $\log e^z = z - 2\pi\I k$ für $-\pi + 2\pi k < \operatorname{Im}(z) < \pi + 2\pi k$.
    \item[iv)] Für $\log(-1)$ gilt der Hauptzweig nicht mehr. Mit Zweigen $\theta \neq -\pi$ gilt: $\log(-1) = \I\pi + 2\pi\I k$ für $k \in \Z$.
\end{itemize}
\end{tcolorbox}

\begin{tcolorbox}[colback=white, colframe=black!70, title={Bemerkung 6.3: Logarithmus-Rechenregeln}]
\begin{itemize}
    \item[i)] Reeller Log $\neq$ komplexer Log.
    \item[ii)] Der Hauptzweig ist nicht für $\log(z^2) = 2\log z$ allgemein definiert, da z.B. $\log(i) = \I\pi/2$, aber $\log(i^2) = \log(-1) = \I\pi \neq 2(\I\pi/2) = \I\pi$ (hier passt es, aber für andere Zweige/Werte nicht).
    \item[iii)] Für $w, z \in \C \setminus (-\infty, 0]$ mit geeigneten Argumenten gilt $\log(wz) = \log w + \log z$.
    \item[iv)] Allgemein gilt $\log(wz) - \log w - \log z \in \{-2\pi\I, 0, 2\pi\I\}$.
\end{itemize}
\end{tcolorbox}

\begin{satz}{6.4 Holomorphie des Hauptzweiges}
Sei $\log$ der Hauptzweig des komplexen Logarithmus auf $\C \setminus (-\infty, 0]$. Dann ist $\log \in H(\C \setminus (-\infty, 0])$ mit $\log'(z) = \frac{1}{z}$. 
Für $|z|<1$ gilt die Potenzreihenentwicklung:
\[
    \log(1+z) = \sum_{k=1}^\infty \frac{(-1)^{k-1}}{k} z^k = z - \frac{z^2}{2} + \frac{z^3}{3} - \dots
\]
\end{satz}

\begin{satz}{6.5 Integral des Logarithmus}
Ist $\gamma: [a,b] \to \C \setminus (-\infty, 0]$ ein $C_{st}^1$-Weg, so gilt:
\[
    \int_\gamma \frac{1}{z} \, \dz = \log \gamma(b) - \log \gamma(a)
\]
Insbesondere ist $\int_\gamma \frac{1}{z} \, \dz = 0$ für jeden geschlossenen Weg mit $Sp(\gamma) \subset \C \setminus (-\infty, 0]$.
\end{satz}

\begin{tcolorbox}[colback=white, colframe=black!70, title={Bemerkungen 6.6 \& 6.7}]
\begin{itemize}
    \item[6.6] Andere Zweige des Logarithmus lassen sich auf den Hauptzweig zurückführen und sind ebenfalls holomorph mit Ableitung $1/z$.
    \item[6.7] Ist $z \in \C \setminus \{0\}$, so heißt jede Zahl $w$ mit $\exp(w) = z$ ein Logarithmus von $z$. Dieser ist eindeutig bis auf Addition von $2\pi\I k, k \in \Z$.
\end{itemize}
\end{tcolorbox}

\begin{defi}{6.8 Komplexe Potenzen}
Sei $a \in \C \setminus \{0\}, b \in \C$ und $\widetilde{\log} a$ ein Logarithmus von $a$. Dann ist
\[
    a^b := \exp(b \widetilde{\log} a)
\]
die (eine) $b$-te Potenz von $a$. Ist $b = 1/n$ ($n \in \N$), so heißt dies \textbf{$n$-te Wurzel}.
\end{defi}

\begin{tcolorbox}[colback=white, colframe=black!70, title={Bemerkung 6.9: Eigenschaften der Potenz}]
\begin{itemize}
    \item[i)] $\widetilde{\log} a$ ist bis auf $2\pi\I k$ definiert $\implies a^b$ eindeutig bis auf Multiplikation mit $\exp(2\pi\I k b)$.
    \item[ii)] Notation $e^b$ meint $\exp(b)$, außer es ist explizit etwas anderes gemeint.
    \item[iii)] $a^{1/n}$ ist die $n$-te Wurzel, eindeutig bis auf Multiplikation mit $\exp(2\pi\I k / n)$.
    \item[iv)] Bei gleichem Logarithmus-Zweig gilt $a^{b_1} a^{b_2} = a^{b_1+b_2}$.
    \item[v)] Beispiel: $i^i = \exp(i \log i) = \exp(i \cdot \I\pi/2) = e^{-\pi/2}$ (reell!).
\end{itemize}
\end{tcolorbox}

\begin{defi}{6.10 Zweig der komplexen Potenz}
Sei $\log$ ein Zweig des Logarithmus. Dann heißt $z \mapsto z^b := \exp(b \log z)$ ein \textbf{Zweig der $b$-ten Potenz}. Wird der Hauptzweig verwendet, spricht man vom Hauptzweig der Potenz.
\end{defi}

\begin{satz}{6.11 Ableitung komplexer Potenzen}
Sei $b \in \C$. Dann ist jeder Zweig der $b$-ten Potenz holomorph und es gilt:
\[
    \frac{d}{dz} (z^b) = b z^{b-1}
\]
(wobei $z^{b-1}$ mit dem gleichen Zweig wie $z^b$ definiert ist).
\end{satz}

% --- Fortführung Kapitel 6: Windungszahl ---

\begin{defi}{6.12 Umlaufzahl (Windungszahl)}
Sei $G \subset \C$ und $\gamma: [a,b] \to G$ ein geschlossener $C_{st}^1$-Weg. Sei $z_0 \notin Sp(\gamma)$. Dann heißt
\[
    n(\gamma, z_0) := \frac{1}{2\pi\I} \int_\gamma \frac{1}{z-z_0} \, \dz
\]
die \textbf{Umlaufzahl} oder \textbf{Windungszahl} des Weges $\gamma$ um $z_0$. Es gilt $n(\gamma, z_0) \in \Z$.
\end{defi}

\begin{tcolorbox}[colback=white, colframe=black!70, title={Intuition zur Windungszahl}]
Die Windungszahl misst, wie oft sich der Weg $\gamma$ um $z_0$ herumwindet (gegen den Uhrzeigersinn positiv).
\begin{itemize}
    \item Der Imaginärteil des Integrals $\int_\gamma \frac{1}{z-z_0} \dz$ entspricht der Zunahme des Arguments entlang des Weges.
    \item Da $\gamma$ geschlossen ist, folgt für die Windungszahl:
    \[
        n(\gamma, z_0) = \frac{1}{2\pi} \cdot \text{Zunahme des Arguments} = \frac{1}{2\pi} \operatorname{Im} \left( \int_\gamma \frac{1}{z-z_0} \, \dz \right)
    \]
\end{itemize}
\end{tcolorbox}

\begin{tcolorbox}[colback=blue!5!white, colframe=blue!75!black, title={Beweisskizze für $n(\gamma, z_0) \in \Z$}]
Betrachte die Funktion $\varphi(t) := \exp \left( \int_a^t \frac{\gamma'(s)}{\gamma(s)-z_0} \, ds \right)$.
\begin{itemize}
    \item Man zeigt, dass $\frac{\varphi(t)}{\gamma(t)-z_0}$ konstant ist, da die Ableitung nach $t$ verschwindet.
    \item Wegen $\gamma(a) = \gamma(b)$ folgt $\varphi(a) = 1$ und $\varphi(b) = \exp \left( \int_\gamma \frac{1}{z-z_0} \, \dz \right)$.
    \item Da $\varphi(b) = \varphi(a) \cdot \frac{\gamma(b)-z_0}{\gamma(a)-z_0} = 1$, muss der Exponent ein Vielfaches von $2\pi\I$ sein, woraus $n(\gamma, z_0) \in \Z$ folgt.
\end{itemize}
\end{tcolorbox}

